\documentclass{beamer}

% --- Pacotes essenciais (seguros com o tema SINTEF) ---
\usepackage[T1]{fontenc}
\usepackage[utf8]{inputenc} % remova se compilar com XeLaTeX/LuaLaTeX
\usepackage[portuguese]{babel}
\usepackage{lmodern}
\usepackage{amsmath,amsfonts}
\usepackage{oldgerm} % você já usava; pode remover se não precisar

% --- Tema e fontes ---
\usetheme{sintef}
\usefonttheme[onlymath]{serif}

% --- Math Bold Font ---
\usepackage{bm} % carregar depois de todas as fontes

% --- Macros do seu template/exemplos ---
\newcommand{\testcolor}[1]{\colorbox{#1}{\textcolor{#1}{test}}~\texttt{#1}}
\newcommand{\hrefcol}[2]{\textcolor{cyan}{\href{#1}{#2}}}

% --- Title background (versão estrelada, antes do \maketitle) ---
\titlebackground*{assets/background}

% --- Bibliografia: mostrar número do item ---
\setbeamertemplate{bibliography item}{\insertbiblabel}

% --- Remove subtítulos de frames ---
\makeatletter
  \def\beamer@checkframetitle{\@ifnextchar\bgroup\beamer@inlineframetitle{}}
\makeatother

% --- Metadados ---
\title{Aprendizado de Representações com Modelos de Linguagem}
\author{Felipe Grael}
\date{\today}

\begin{document}

\maketitle

\begin{frame}{Contexto}

	\begin{columns}[t]

		\column{0.5\textwidth}
		\vspace{-5mm}

    \begin{center}
    \textbf{Objetivo}
    \end{center}

    Derivar \textbf{representações vetoriais} de documentos respeitando relações semânticas.

    \small
    \begin{itemize}
      \item Similaridade: documentos similares devem ter vetores próximos
      \item Recuperação: obter documentos relevantes para um contexto
      \item \textbf{Clusterização / Tópicos}: descobrir agrupamentos e assuntos em comum entre documentos
    \end{itemize}

		\column{0.5\textwidth}
		\vspace{-5mm}

    \begin{center}
    \textbf{Histórico}
    \end{center}
    \small
    \begin{itemize}
      \item Vector Space Model e TF-IDF trabalham somente com presença de palavras.
      \item \emph{Word Embeddings} começam a capturar valor semântico de palavras, mas sem contexto.
      \item Modelos de Linguagem modernos, especialmente os Grandes (LLMs), se mostram promissores em capturar nuanças semânticas e contextuais.
    \end{itemize}

  \end{columns}

\end{frame}


\begin{frame}{Problema}

	\begin{columns}[t]
		\column{0.5\textwidth}
		\vspace{-5mm}

    \begin{center}
    \textbf{Problema}
    \end{center}

    \small
    \begin{itemize}
      \item Modelos de linguagem geram boas representações para tokens, mas não para frases ou documentos.
      \item Pode-se usar tokens especiais ou médias mas com eficácia limitada.
      \item \textbf{Anisotropia} dos espaços de embedding: vetores tendem a se concentrar em subespaços estreitos.

        \begin{itemize}
          \item \textbf{Todos os vetores possuem similaridade alta!}
        \end{itemize}

    \end{itemize}

		\column{0.5\textwidth}
		\vspace{-5mm}

    \begin{center}
    \textbf{Trabalhos Relacionados}
    \end{center}

    \small
    \begin{itemize}
      \item \textbf{Sentence BERT}\cite{reimers2019sentencebertsentenceembeddingsusing} - fine-tuning com "redes siamesas" para aprender similaridade entre sentenças.
      \item \textbf{Contrastive Learning} - fine tuning com função custo que incentiva a proximidade de sentenças similares e distância entre sentenças diferentes.
      \item Desafios em dataset de qualidade para fine-tuning, e abordagens de aumento de dados (\emph{data augmentation})
    \end{itemize}

	\end{columns}
\end{frame}


\begin{frame}{Características de LLM}

  \begin{center}
    \large\texttt{[BOS] I want to drink some orange \textbf{[MASK]} . [EOS]}
  \end{center}

  Preenchimento de máscara:

  \begin{center}
    \begin{tabular}{c|l}
      Token & logits \\
      \hline
      \texttt{juice} & 14.73 \\
      \texttt{\#\#s} & 10.829 \\
      \texttt{wine} & 9.72 \\
      \texttt{soda} & 9.25 \\
      \texttt{tea} & 8.49
    \end{tabular}
  \end{center}

\end{frame}


\begin{frame}{Scaling Sentence Embedings with LLM \cite{jiang2024prompteol}\footnotemark}

  \begin{center}
    \large \texttt{This sentence: `` [text] '' means in one word `` [MASK] '' .}
  \end{center}

  \begin{itemize}
    \item \textbf{PromptEOL} - \emph{Prompt-based method with Explicit One word Limitation}
    \item Explora último estado interno antes da camada de saída do modelo como representação
    \item Bons resultados em modelos abertos (OPT)
    \item Trabalho explora few-shot e fine tuning também.
  \end{itemize}

  \footnotetext[1]{Findings of the Association for Computational Linguistics: EMNLP 2024}

\end{frame}


\begin{frame}{Método}

  \textbf{Objetivo}: Avaliar clusterização por embeddings obtidos por PromptEOL.

	\begin{columns}[t]

		\column{0.5\textwidth}

    \small
    \begin{itemize}
      \item Estudar dataset \textbf{20 Newsgroups}
      \item Obter embeddings em diferentes modelos
      \item Clusterização com K-means (+~bootstrapping)
      \item Avaliação:

      \begin{itemize}
        \item Intrínseca: Silhouette
        \item Extrínseca: Adjusted Rand Index (ARI)
      \end{itemize}

      \item Baseline: K-Means com TF-IDF
    \end{itemize}

		\column{0.5\textwidth}

    Modelos:
    \small
    \begin{itemize}
      \item BERT
      \item BERT + KAN
      \item LLaDA (Discrete Diffusion)
    \end{itemize}

    Somente com pré-treino, executando tarefa MLM (Masked Language Modeling).

  \end{columns}

\end{frame}


\begin{frame}{20 Newsgroups}

	\begin{columns}[t]

		\column{0.7\textwidth}
    \small
    \begin{itemize}
      \item 18331 documentos
      \item Pré-processados para remover cabeçalhos
      \item Truncados para o número máximo de tokens do modelo
    \end{itemize}

		\column{0.3\textwidth}
		\vspace{-5mm}
    \tiny
    \begin{itemize}
      \item alt.atheism
      \item comp.graphics
      \item comp.os.ms-windows.misc
      \item comp.sys.ibm.pc.hardware
      \item comp.sys.mac.hardware
      \item comp.windows.x
      \item misc.forsale
      \item rec.autos
      \item rec.motorcycles
      \item rec.sport.baseball
      \item rec.sport.hockey
      \item sci.crypt
      \item sci.electronics
      \item sci.med
      \item sci.space
      \item soc.religion.christian
      \item talk.politics.guns
      \item talk.politics.mideast
      \item talk.politics.misc
      \item talk.religion.misc
    \end{itemize}

  \end{columns}

\end{frame}

\begin{frame}{Pipeline}
  \begin{figure}[h]
    \centering
    \includegraphics[height=1.1\textheight]{assets/pipeline.png}
  \end{figure}
\end{frame}

\begin{frame}{Resultados}
  \centering
  \tiny
  \begin{tabular}{ l | c | c | c | c }
    Modelo & Arquitetura & Embedding & Silhouette & ARI \\
    \hline
    TF-IDF                     & -                       & $5000$   & $ 0.0097 \  \sigma\!=\!0.0005 $ & $ 0.085 \  \sigma = 0.0078 $ \\
    bert-base-uncased (OPT)    & $ l=12, h=12, 109.5M $ & $768$    & $ 0.049  \  \sigma\!=\!0.002  $ & $ \mathbf{0.26} \  \sigma = 0.006 $ \\
    bert-small                 & $ l=6, h=6, 22.5M $    & $384$    & $ \mathbf{0.10}   \  \sigma\!=\!0.002  $ & $ 0.01 \  \sigma = 0.0006 $ \\
    bert-smaller               & $ l=6, h=4, 8.6M $     & $192$    & $ 0.024  \  \sigma\!=\!0.002  $ & $ 0.06 \  \sigma = 0.002  $ \\
    modernbert-1.5b            & $ l=6, h=4, 8.6M $     & $192$    & $ 0.013  \  \sigma\!=\!0.001  $ & $ 0.05 \  \sigma = 0.0003  $ \\
    modernbert-kan-1.5b        & $ l=6, h=4, 8.6M $     & $192$    & $ 0.011  \  \sigma\!=\!0.001  $ & $ 0.004 \  \sigma = 0.0003  $ \\
    modernbert-diffusion-1.5b  & $ l=6, h=4, 8.6M $     & $192$    & $ 0.012  \  \sigma\!=\!0.001  $ & $ 0.005 \  \sigma = 0.0003  $ \\
    LLaDA-8B-Base              & $ l=32, h=32, 8B $     & $126336$ &  ? & ?
  \end{tabular}

\end{frame}

\begin{frame}{Resultados - bert-small}

  \begin{center}
    Projecão t-SNE dos embeddings obtidos com \texttt{bert-small}, com melhor Silhouette.
  \end{center}

  \begin{figure}[h]
    \centering
    \includegraphics[height=\textheight]{assets/bert-small-tsne.png}
  \end{figure}


\end{frame}

\begin{frame}{Conclusões}

  \begin{itemize}
    \item Resultados tendem a melhorar com mais dados de pré-treino
    \item Possibilidade de controlar tokenização com modelos pré-treinados com volumes médios de dados
    \item Método promissor para exploração de relações semânticas com nuançadas
  \end{itemize}

\end{frame}

%===========================
% REFERÊNCIAS
%===========================
\begin{frame}[allowframebreaks]
	\frametitle{Referências Bibliográficas}
	\bibliographystyle{ieeetr}
	\bibliography{presentation_bib.bib}
\end{frame}

\end{document}
